\documentclass[10pt,a4paper]{article}
\usepackage[utf8]{inputenc}
\usepackage[margin=0.7in]{geometry}
\usepackage{xcolor}
\usepackage{titlesec}
\usepackage{enumitem}
\usepackage[hidelinks]{hyperref}
\hypersetup{pdftitle={Varun Ravi Varma - CV (AI Research Scientist / Research Engineer)}, pdfauthor={Varun Ravi Varma}}

\pagestyle{empty}
\raggedbottom
\raggedright
\setlength{\tabcolsep}{0in}

% Section formatting
\titleformat{\section}{
  \vspace{-3pt}\scshape\raggedright\large
}{}{0em}{}[\color{black}\titlerule \vspace{-4pt}]

\begin{document}
% Public web version of the industry CV, published on v3rm1.github.io.
% Personal details (address, phone, date of birth, citizenship, visa status)
% and unpublished research specifics are deliberately left out.
% Build: pdflatex industry-cv.tex, then copy the PDF to ../files/cv/.

% Heading / Personal Details
\begin{center}
    {\Huge \scshape Varun Ravi Varma} \\ \vspace{2pt}
    {\large AI Research Scientist \& AI / Research Engineer} \\ \vspace{3pt}
    \small Grimstad, Norway $|$ \href{mailto:varunravivarma@gmail.com}{varunravivarma@gmail.com} $|$ \href{https://v3rm1.github.io/}{v3rm1.github.io} \\
    GitHub: \href{https://github.com/v3rm1}{v3rm1} $|$ ORCID: \href{https://orcid.org/0000-0003-1719-2478}{0000-0003-1719-2478} $|$ LinkedIn: \href{https://www.linkedin.com/in/varun-ravi-varma}{varun-ravi-varma}
\end{center}

% Availability & languages
\section{Availability \& Languages}
\begin{itemize}[leftmargin=0.15in, label={}]
    \small{\item{
     \textbf{Availability}{: Research programmer and applied scientist roles in Germany, Austria or Switzerland (DACH) from early 2028} \\
     \textbf{Languages}{: English (C2 -- Proficient), German (A2--B1, currently taking courses at Heidelberg University)}
    }}
\end{itemize}

% Professional Summary
\section{Professional Summary}
AI Research Scientist and Software Engineer with over 5 years of experience bridging theoretical machine learning breakthroughs and production-grade software architectures.
Specialised in GPU-accelerated simulation platforms, multi-agent reinforcement learning (MARL), LLM/agentic workflows, and semantic knowledge graph pipelines.
Demonstrated track record of transforming complex mathematical formalisms into high-throughput, reproducible research software across distributed GPU clusters (HPC) and enterprise cloud environments.
Experienced in designing robust data pipelines, establishing CI/CD automation, and deploying resilient machine learning systems in healthcare, finance, and industrial domains.

% Technical Skills
\section{Technical Skills}
\begin{itemize}[leftmargin=0.15in, label={}]
    \small{\item{
     \textbf{Programming Languages}{: \textbf{Python}, \textbf{C++}, \textbf{Rust}, \textbf{SQL}, \textbf{CUDA}, \textbf{Bash}, \textbf{PowerShell}} \\
     \textbf{AI Frameworks \& Modelling}{: \textbf{PyTorch}, \textbf{JAX}, \textbf{Multi-Agent Systems (MARL)}, \textbf{LLMs \& Agentic Architectures}, \textbf{Graph Neural Networks}, \textbf{Scientific ML}, \textbf{Non-Linear Dynamical Systems}} \\
     \textbf{Systems, Cloud \& MLOps}{: \textbf{Docker}, \textbf{Kubernetes}, \textbf{Git-based CI/CD}, \textbf{Linux Server Administration}, \textbf{Distributed GPU Clusters (Slurm/HPC)}, \textbf{Azure}, \textbf{Databricks}, \textbf{Model Serving (REST/gRPC)}}
    }}
\end{itemize}

% Professional Experience
\section{Professional Experience}
\textbf{Doctoral Research Fellow (AI \& Complex Systems)} \hfill 09/2023 -- Present \\
\textit{University of Agder} \hfill \textit{Grimstad, Norway}
\begin{itemize}[leftmargin=0.15in, label=$\bullet$]
    \item Research and implement multi-agent reinforcement learning (MARL) algorithms for coordinated, long-horizon decision-making in complex multi-agent systems.
    \item Architected a high-throughput, GPU-vectorised simulation engine in JAX (\texttt{jit}, \texttt{vmap}) that steps thousands of scenarios in parallel on a single GPU.
    \item Designed an automated validation suite to check simulated dynamics, parameter sensitivity, and policy behaviour, with reproducible multi-seed experiment sweeps on multi-GPU servers.
    \item Mentored student researchers in modular software design, version control workflows, unit testing, and reproducible code standards for scientific machine learning.
    \item Directed advanced scientific writing and technical presentation workshops across university departments.
\end{itemize}

\vspace{3pt}
\textbf{Scientific Programmer} \hfill 06/2022 -- 08/2023 \\
\textit{Sano Centre for Computational Medicine} \hfill \textit{Krakow, Poland}
\begin{itemize}[leftmargin=0.15in, label=$\bullet$]
    \item Architected analytical data pipelines and knowledge graph reasoning frameworks for clinical decision support and complex biomedical artifact analysis.
    \item Designed relational schemas and semantic dependency mapping algorithms to structure heterogeneous medical data into queryable knowledge graphs.
    \item Developed and containerised reproducible computational tools and model evaluation harnesses using Docker to ensure seamless cross-platform execution.
    \item Implemented Git-based CI/CD workflows to automate regression testing, model training pipelines, and data integrity verification in regulated environments.
    \item Collaborated with clinicians and domain scientists to translate rigorous clinical requirements into production-grade software architectures.
\end{itemize}

\vspace{3pt}
\textbf{Teaching Assistant} \hfill 08/2025 -- 01/2026 \\
\textit{University of Agder} \hfill \textit{Grimstad, Norway}
\begin{itemize}[leftmargin=0.15in, label=$\bullet$]
    \item Research Methods for AI and IoT: assisted with lectures, tutorials and exam evaluations.
    \item Computer Vision: assisted with lectures and labs, mentored student projects and evaluated technical reports.
\end{itemize}

\vspace{3pt}
\textbf{Graduate Teaching Assistant} \hfill 02/2020 -- 08/2020 \\
\textit{University of Groningen} \hfill \textit{Groningen, Netherlands}
\begin{itemize}[leftmargin=0.15in, label=$\bullet$]
    \item Deep Learning: assisted with labs, mentored students and evaluated technical practical reports.
    \item Logical Aspects of Multi-Agent Systems: facilitated tutorials, graded assignments and reviewed course projects.
\end{itemize}

\vspace{3pt}
\textbf{Consultant -- Data Analytics \& Machine Learning} \hfill 06/2016 -- 11/2018 \\
\textit{Neudesic Global Services} \hfill \textit{Bangalore, India}
\begin{itemize}[leftmargin=0.15in, label=$\bullet$]
    \item Architected and deployed enterprise-scale data analytics pipelines, deep learning models, and high-capacity data warehousing solutions for Fortune 500 clients.
    \item Developed proof-of-concept machine learning solutions across energy, pharmaceutical, and financial sectors.
    \item Managed migration and integration pipelines for structured and unstructured data across hybrid cloud and on-premise high-availability infrastructures.
    \item Founded and led an internal technical upskilling initiative, directing weekly seminars on statistics, machine learning, and scalable systems.
    \item Automated maintenance procedures and system monitoring via scripting to ensure 99.9\% operational uptime.
\end{itemize}

% Selected Research & Engineering Projects
\section{Key Research \& Engineering Projects}
\textbf{GPU-Accelerated Multi-Agent Simulation} $|$ \textit{JAX, Python} $|$ \textit{Ongoing, details after publication} \\
\begin{itemize}[leftmargin=0.15in, label=$\bullet$]
    \item Building fully vectorised, JIT-compiled simulation environments in JAX so that thousands of scenarios step in parallel on a single GPU.
    \item Validation suite to check that the simulated dynamics behave plausibly.
\end{itemize}

\vspace{2pt}
\textbf{Hierarchical Multi-Agent Reinforcement Learning} $|$ \textit{JAX, PyTorch, Python} $|$ \textit{Ongoing, details after publication} \\
\begin{itemize}[leftmargin=0.15in, label=$\bullet$]
    \item Designing structured, hierarchical MARL architectures in which agents coordinate in sub-teams under partial observability and long horizons.
    \item Compared against established flat MARL baselines across many random seeds with statistically sound aggregate metrics.
\end{itemize}

\vspace{2pt}
\textbf{Clinical Knowledge Graph \& Semantic Reasoning Engine} $|$ \textit{Python, Neo4j, Docker, CI/CD} \\
\begin{itemize}[leftmargin=0.15in, label=$\bullet$]
    \item Built end-to-end knowledge extraction and semantic reasoning graph pipelines unifying structured electronic health records and unstructured medical literature.
    \item Containerised pipeline components with Docker and integrated automated validation testing into continuous integration workflows.
\end{itemize}

% Education
\section{Education}
\textbf{Ph.D. in Artificial Intelligence} \hfill 2023 -- Present \\
\textit{University of Agder} \hfill \textit{Grimstad, Norway} \\
\textit{Focus: Multi-Agent Reinforcement Learning, GPU-Accelerated Simulations, and Complex Systems}
\vspace{4pt}

\textbf{M.Sc. in Artificial Intelligence} \hfill 2019 -- 2021 \\
\textit{University of Groningen} \hfill \textit{Groningen, Netherlands}
\vspace{4pt}

\textbf{B.Tech. in Mechanical Engineering} \hfill 2012 -- 2016 \\
\textit{Visveswaraya Technological University} \hfill \textit{Belagavi, India}

\end{document}
